%%%%%%%%%%%%%%%%%%%%%%%%%%%%%%%%%%%%%%%%%%%%%%%%%%%%%%%%%%%%%%%%
%
%  Template for creating scribe notes.
%
%  Fill in your name, lecture number, lecture date and body
%  of scribe notes as indicated below.
%
%%%%%%%%%%%%%%%%%%%%%%%%%%%%%%%%%%%%%%%%%%%%%%%%%%%%%%%%%%%%%%%%


\documentclass[11pt]{article}

\usepackage{graphicx}
\usepackage{amssymb}

\setlength{\topmargin}{0pt}
\setlength{\textheight}{9in}
\setlength{\headheight}{0pt}
\setlength{\headsep}{0pt}
\setlength{\oddsidemargin}{0.25in}
\setlength{\textwidth}{6in}

\newcommand{\draftnotice}{\vbox to 0.25in{\noindent
   \raisebox{0.6in}[0in][0in]{\makebox[\textwidth][r]{\it
    DRAFT --- a final version will be posted shortly}}}
   \vspace{-.25in}\vspace{-\baselineskip}
}

\pagestyle{plain}

\begin{document}

\thispagestyle{empty}


\begin{center}
\bf\large Introduction to Data Structure and Algorithm
\end{center}

\noindent
Lecturer: Jalaj Upadhyay   %%% FILL IN LECTURER (if not RS)
\hfill
Lecture 1\               %%% FILL IN LECTURE NUMBER HERE
\\
Scribe: HaoyuWang                 %%% FILL IN YOUR NAME HERE
\hfill
Jan 18th, 2022           %%% FILL IN LECTURE DATE HERE

\noindent
\rule{\textwidth}{1pt}

\medskip

%%%%%%%%%%%%%%%%%%%%%%%%%%%%%%%%%%%%%%%%%%%%%%%%%%%%%%%%%%%%%%%%
%% BODY OF SCRIBE NOTES GOES HERE
%%%%%%%%%%%%%%%%%%%%%%%%%%%%%%%%%%%%%%%%%%%%%%%%%%%%%%%%%%%%%%%%
The Algothrim is sequence of atomic steps to tell a machine how to solve a problem. 
Any ways to solve the problem can be called a algorithm. 
But the main purpose is to find the balance between the memory usage and the solving time of current target problem.

{\bf Algorithm 1: Minimum Finding}

Just go through the whole Array $A$ of size $n$ numbers to find the current smallest element by setting a pointer to the latest smallest element and a pointer to move from $A[0]$ to $A[n]$
The running time denpended on the size of Array $A$ because we need $n$ atomic steps to go through and compare to find the smallest element.













%%%%%%%%%%%%%%%%%%%%%%%%%%%%%%%%%%%%%%%%%%%%%%%%%%%%%%%%%%%%%%%%

\end{document}
